\chapter{Resultados}
\label{cap:04}

\section{Apresentação dos resultados}

\subsection{Experimento 1 — Classificação \textit{Zero-Shot}}

O CLIP ViT-L/14 em modo \textit{zero-shot} obteve acurácia de 53,85\% e AUC-ROC de 0,5392 sobre o conjunto \texttt{test\_unseen} (2.000 amostras). Para a classe de ódio, a precisão foi de 0,41, a revocação de 0,50 e o F1-Score de 0,45; o F1 macro entre as classes ficou em 0,53, apenas marginalmente acima do acaso.

A aplicação do \textbf{Classificador Fuzzy Unificado} sobre as probabilidades do CLIP resultou em acurácia de 59,30\%. O sistema tornou-se mais conservador: a precisão para a classe de ódio subiu para 0,40, mas a revocação caiu para 0,17 (F1-Score de 0,24). A distribuição das categorias revelou que 62,4\% das amostras foram alocadas em \textit{Ambiguidade/Moderado} (1.247 de 2.000), reflexo direto da alta entropia das predições zero-shot.

O estudo de ablação por modalidade evidenciou um resultado contraintuitivo: a abordagem \textbf{Somente Texto} alcançou a maior acurácia (61,80\%), superando tanto \textbf{Somente Imagem} (53,05\%) quanto o modo \textbf{Multimodal} (53,85\%). Isso indica que, sem ajuste fino, o CLIP não consegue capturar a justaposição contraditória entre imagem e texto que caracteriza os memes de ódio.

\subsection{Experimento 2 — Pré-processamento com Agrupamento Semântico}

O \texttt{MultimodalModel} com enriquecimento por clusters obteve acurácia de 67,55\%, AUC-ROC de 0,7318 e \textit{Average Precision} de 0,6074 em \texttt{test\_unseen}. O modelo apresentou tendência conservadora: precisão de 0,62 e revocação de 0,36 para a classe de ódio (F1 de 0,45); F1 macro de 0,61.

Com a \textbf{Lógica Fuzzy}, a acurácia binarizada manteve-se em 66,60\%, mas o comportamento mudou significativamente, a precisão para ódio subiu para 0,72, enquanto a revocação caiu para 0,18 (F1 de 0,29). O sistema deslocou a maioria das amostras para \textit{Nenhum Ódio} (758) e \textit{Ambiguidade/Moderado} (562), classificando apenas 188 exemplos como ódio efetivo (\textit{Ódio}: 146; \textit{Alto Grau}: 42).

O estudo de ablação com cinco configurações demonstrou a contribuição incremental de cada componente, conforme a Tabela~\ref{tab:ablacao_exp2}.

\begin{table}[htbp]
\centering
\caption{Estudo de ablação — Experimento 2 (\texttt{test\_unseen}, 2.000 amostras)}
\label{tab:ablacao_exp2}
\begin{tabular}{lcccc}
\hline
\textbf{Configuração} & \textbf{Acurácia} & \textbf{AUC-ROC} & \textbf{F1 macro} & \textbf{F1 ódio} \\
\hline
Somente Imagem    & 63,60\% & 0,7188 & 0,44 & 0,10 \\
Somente Texto     & 63,25\% & 0,6310 & 0,48 & 0,20 \\
CLIP Multimodal   & 67,60\% & 0,7321 & 0,66 & 0,58 \\
+ Cluster Texto   & 67,90\% & 0,7318 & 0,65 & 0,49 \\
Modelo Completo   & \textbf{68,65\%} & \textbf{0,7336} & 0,63 & 0,49 \\
\hline
\end{tabular}
\end{table}

\subsection{Experimento 3 — Pipeline Completo com HateCLIPper}

O HateCLIPper com fusão \textit{cross} (FIM) e limiar ótimo de 0,0119 (determinado por maximização do F1 em \texttt{dev\_seen}) obteve \textbf{AUROC de 0,8121}, \textbf{Acurácia de 75,55\%} e \textbf{F1-Score de 0,6407} em \texttt{test\_unseen}. O modelo demonstrou alta sensibilidade para a classe de ódio (revocação de 0,86, precisão de 0,55) e boa especificidade para não ódio (precisão de 0,87, revocação de 0,58), refletindo o ajuste do limiar para maximizar F1 em detrimento da precisão.

O estudo de ablação entre os quatro modos de fusão confirma a superioridade do modo \textit{cross}, conforme a Tabela~\ref{tab:ablacao_exp3}.

\begin{table}[htbp]
\centering
\caption{Estudo de ablação — modos de fusão do HateCLIPper (\texttt{test\_unseen})}
\label{tab:ablacao_exp3}
\begin{tabular}{lcccc}
\hline
\textbf{Fusão} & \textbf{Val AUROC} & \textbf{Test AUROC} & \textbf{Acurácia} & \textbf{F1 macro} \\
\hline
align        & 0,7944 & 0,8156 & 75,40\% & 0,6227 \\
concat       & 0,7536 & 0,7778 & 72,10\% & 0,5798 \\
\textbf{cross} & \textbf{0,7962} & \textbf{0,8141} & \textbf{75,90\%} & \textbf{0,6456} \\
align\_concat & 0,7529 & 0,7798 & — & — \\
\hline
\end{tabular}
\end{table}

O \textbf{Sistema Fuzzy Mamdani} aplicado sobre os \textit{scores} visuais e textuais independentes do HateCLIPper produziu um limiar ótimo de 7,17 (escala de 0 a 100). A distribuição sobre \texttt{test\_unseen} resultou em: \textit{Inofensivo} 1.063 (53,2\%), \textit{Questionável} 300 (15,0\%), \textit{Moderado} 279 (14,0\%), \textit{Crítico} 217 (10,9\%) e \textit{Extremo} 141 (7,1\%).

\section{Análise crítica e comparação}

\subsection{Progressão entre experimentos}

A evolução progressiva dos três experimentos evidencia ganhos consistentes: a acurácia passou de 53,85\% (\textit{zero-shot}) para 67,55\% (pré-processamento com clusters) e 75,55\% (\textit{fine-tuning} com HateCLIPper), com o AUROC saltando de 0,5392 para 0,7318 e 0,8121, respectivamente. O resultado do Experimento~3 aproxima-se do nível humano reportado no \textit{Hateful Memes Challenge} (AUROC $= 0{,}820$), superando os \textit{baselines} do VisualBERT e do CLIP \textit{zero-shot} da literatura (ambos com AUROC $\approx 0{,}730$), embora ainda aquém do HateCLIPper original com \textit{cross-validation} complexa (AUROC $= 0{,}859$).

\subsection{Papel da lógica \textit{fuzzy}}

Em todos os experimentos, a lógica \textit{fuzzy} introduziu um \textit{trade-off} consistente: maior precisão para ódio (redução de falsos positivos) em troca de menor revocação (aumento de falsos negativos). No Experimento~1, a revocação caiu de 0,50 para 0,17; no Experimento~2, de 0,36 para 0,18. Esse comportamento é desejável em cenários de moderação real, onde classificar erroneamente conteúdo benigno como ódio tem custo alto, mas impõe a necessidade de revisão humana para os casos enquadrados como \textit{Ambiguidade/Moderado}. No Experimento~3, o sistema Mamdani com dois antecedentes independentes (score visual e textual) produz uma graduação mais informativa do que a abordagem baseada em entropia dos experimentos anteriores, pois captura assimetrias entre as modalidades.

\subsection{Limitações}

O principal gargalo em todos os experimentos é a baixa revocação para a classe de ódio, especialmente após a aplicação da lógica \textit{fuzzy}. A dificuldade em capturar sarcasmo, ironia e referências socioculturais implícitas não é resolvida apenas pela fusão de modalidades, exigindo formas de raciocínio de senso comum que os modelos atuais não incorporam plenamente.

\section{Discussão}

Os resultados confirmam a hipótese central de que a fusão multimodal supervisionada supera abordagens \textit{zero-shot} e unimodais na detecção de discurso de ódio em memes. O fato de o modo \textit{Somente Texto} superar o \textit{Multimodal} no Experimento~1 e de as abordagens unimodais do Experimento~2 ficarem abaixo do CLIP Multimodal, demonstra que o benefício da multimodalidade só se manifesta com treinamento supervisionado adequado.

A progressão dos experimentos cumpre os objetivos específicos desta pesquisa: estabelecer uma linha de base (\textit{zero-shot}), investigar o enriquecimento por agrupamento semântico (pré-processamento) e implementar um \textit{fine-tuning} completo com fusão cruzada (HateCLIPper). A integração da lógica \textit{fuzzy} em todos os estágios acrescenta interpretabilidade ao sistema, convertendo probabilidades brutas em categorias linguísticas que refletem a gradação inerente ao discurso de ódio.

%Trabalhos futuros devem explorar modelos de fundação multimodais mais recentes, a incorporação de raciocínio de senso comum (\textit{commonsense knowledge}) na etapa de fusão, e o refinamento das regras \textit{fuzzy} com base em características regionais das imagens, buscando aproximar o desempenho ao nível de compreensão humana.
